%! Author = lili
%! Date = 6/15/26

% Seite 23 - ??


\skript{The unit cost edit distance simply counts certain differences between strings.
There is no inherent
biological meaning behind this scheme.
If we want to compare DNA or protein sequences, for example,
    biologically more meaningful distances should be used.}

\defin{transversion_transition_cost}

\skript{The following is called the \emph{transversion/transition cost} on the DNA alphabet:
\[
    \begin{array}{c|cccc}
        & \mathtt{A} & \mathtt{G} & \mathtt{C} & \mathtt{T}\\
        \hline
        \mathtt{A} & 0 & 1 & 2 & 2\\
        \mathtt{G} & 1 & 0 & 2 & 2\\
        \mathtt{C} & 2 & 2 & 0 & 1\\
        \mathtt{T} & 2 & 2 & 1 & 0
    \end{array}
\]
Bases A and G are called \emph{purines}, and bases C and T are called \emph{pyrimidines}.
The
transversion/transition cost function reflects the biological fact that a purine/purine and a
pyrimidine/pyrimidine replacement is more likely to occur than a purine/pyrimidine replacement.
Often, an insertion/deletion cost of $3$ is used with the above costs to take into account that a
deletion or an insertion of a base is even more seldom.
These costs define the
transition/transversion distance between two DNA sequences.

To compare protein sequences and define a cost function on the amino acid alphabet, we may count how
many DNA bases must change minimally in a codon to transform one amino acid into another.

A more sophisticated cost function for replacement of amino acids was calculated by W.\ Taylor
according to the method described in Taylor and Jones (1993) and is shown in Table~4.1. (To
completely define a distance function, we would also have to define the costs for insertions and
deletions.)}

\begin{notte}
  \skript{Note, however, that in a distance model a match has always zero cost,
      $\mathrm{cost}\!\left(\begin{smallmatrix}
                                a\\a
      \end{smallmatrix}\right)=0$ for all $a\in\Sigma$.
      This
      restricts flexibility a little bit because any match is scored in the same way.}
\end{notte}


\begin{mytable}
    \begin{tabular}{c|cccccccccccccccccccc}
        & A  & R  & N  & D  & C  & Q  & E  & G  & H  & I  & L  & K  & M  & F  & P  & S  & T  & W  & Y  & V  \\
        \hline
        A & 0  & 14 & 7  & 9  & 20 & 9  & 8  & 7  & 12 & 13 & 17 & 11 & 14 & 24 & 7  & 6  & 4  & 32 & 23 & 11 \\
        R & 14 & 0  & 11 & 15 & 26 & 11 & 14 & 17 & 11 & 18 & 19 & 7  & 16 & 25 & 13 & 12 & 13 & 25 & 24 & 18 \\
        N & 7  & 11 & 0  & 6  & 23 & 5  & 6  & 10 & 7  & 16 & 19 & 7  & 16 & 25 & 10 & 7  & 7  & 30 & 23 & 15 \\
        D & 9  & 15 & 6  & 0  & 25 & 6  & 3  & 9  & 11 & 19 & 22 & 10 & 19 & 29 & 11 & 11 & 10 & 34 & 27 & 17 \\
        C & 20 & 26 & 23 & 25 & 0  & 26 & 25 & 21 & 25 & 22 & 26 & 25 & 26 & 26 & 22 & 20 & 21 & 34 & 22 & 21 \\
        Q & 9  & 11 & 5  & 6  & 26 & 0  & 5  & 12 & 7  & 17 & 20 & 8  & 16 & 27 & 10 & 11 & 10 & 32 & 26 & 16 \\
        E & 8  & 14 & 6  & 3  & 25 & 5  & 0  & 9  & 10 & 17 & 21 & 10 & 17 & 28 & 11 & 10 & 9  & 34 & 26 & 16 \\
        G & 7  & 17 & 10 & 9  & 21 & 12 & 9  & 0  & 15 & 17 & 21 & 13 & 18 & 27 & 10 & 9  & 9  & 33 & 26 & 15 \\
        H & 12 & 11 & 7  & 11 & 25 & 7  & 10 & 15 & 0  & 17 & 19 & 10 & 17 & 24 & 13 & 11 & 11 & 30 & 21 & 16 \\
        I & 13 & 18 & 16 & 19 & 22 & 17 & 17 & 17 & 17 & 0  & 9  & 17 & 8  & 17 & 16 & 14 & 12 & 31 & 18 & 4  \\
        L & 17 & 19 & 19 & 22 & 26 & 20 & 21 & 21 & 19 & 9  & 0  & 19 & 7  & 14 & 20 & 19 & 17 & 27 & 18 & 10 \\
        K & 11 & 7  & 7  & 10 & 25 & 8  & 10 & 13 & 10 & 17 & 19 & 0  & 15 & 26 & 11 & 10 & 10 & 29 & 25 & 16 \\
        M & 14 & 16 & 16 & 19 & 26 & 16 & 17 & 18 & 17 & 8  & 7  & 15 & 0  & 18 & 17 & 16 & 13 & 29 & 20 & 8  \\
        F & 24 & 25 & 25 & 29 & 26 & 27 & 28 & 27 & 24 & 17 & 14 & 26 & 18 & 0  & 27 & 24 & 23 & 24 & 8  & 19 \\
        P & 7  & 13 & 10 & 11 & 22 & 10 & 11 & 10 & 13 & 16 & 20 & 11 & 17 & 27 & 0  & 9  & 8  & 32 & 26 & 14 \\
        S & 6  & 12 & 7  & 11 & 20 & 11 & 10 & 9  & 11 & 14 & 19 & 10 & 16 & 24 & 9  & 0  & 5  & 29 & 22 & 13 \\
        T & 4  & 13 & 7  & 10 & 21 & 10 & 9  & 9  & 11 & 12 & 17 & 10 & 13 & 23 & 8  & 5  & 0  & 31 & 22 & 10 \\
        W & 32 & 25 & 30 & 34 & 34 & 32 & 34 & 33 & 30 & 31 & 27 & 29 & 29 & 24 & 32 & 29 & 31 & 0  & 25 & 32 \\
        Y & 23 & 24 & 23 & 27 & 22 & 26 & 26 & 26 & 21 & 18 & 18 & 25 & 20 & 8  & 26 & 22 & 22 & 25 & 0  & 20 \\
        V & 11 & 18 & 15 & 17 & 21 & 16 & 16 & 15 & 16 & 4  & 10 & 16 & 8  & 19 & 14 & 13 & 10 & 32 & 20 & 0  \\
    \end{tabular}
    \caption{Amino acid replacement costs as suggested by W.\ Taylor.}
\end{mytable}


\section{Sensible Similarity Scores}\label{sec:sensible-similarity-scores}


\skript{The concept of a metric is useful because it embodies the essential properties of our intuition
about distances.
For example, if the triangle inequality is violated, we could find a shorter way
from $x$ to $y$ via a detour over $z$, which is not intuitive.
However, distances also have
disadvantages, as we have seen in the discussion of local alignment, where distance functions are
not useful.}

\defin{general_similarity_score}

\skript{Similarity scores, instead, where matches are scored positively, while mismatches and
indels receive negative scores, provide a useful alternative.
Here we consider a general similarity
score that should satisfy certain intuitive properties.}

\definSB{sensible_score_function}

\skript{The first and second conditions impose symmetry and state that the similarity between $a$ and itself
is always higher than between $a$ and anything else.
In particular, insertions and deletions never
score positively.
The third condition states that a direct substitution is preferable to a deletion
followed by an insertion.

In general, it is a good idea to verify that a score function is sensible, like in the case of
log-odds score matrices introduced in the next section.
In other cases, one might deviate from this
property for good reason, see Section~\ref{sec:non-symmetric-score-matrices}.}

\section{Log-Odds Score Matrices}\label{sec:log-odds-score-matrices}

\begin{idee}
    Scores sind keine Willkür, sondern Log-Odds: Wie viel wahrscheinlicher ist ein Alignment-Paar
    $(a,b)$ unter dem \emph{Homologie-Modell} $M$ als unter dem Zufalls-Modell $R$?
    Positiver Score $=$ eher homolog.
\end{idee}

\skript{Now it is time to discuss how we can define similarity scores for biological sequences from an
evolutionary point of view; we focus on proteins and define a similarity function on the alphabet
    $\Sigma$ of $20$ amino acids.

    The basic observation is that over evolutionary time-scales, in functional proteins two similar
    amino acids are replaced more frequently by each other than dissimilar amino acids.
    The twist is now
    to define similarity by observing how often one amino acid is replaced by another one in a given
    amount of time.}

\definSB{frequency_vector}

\skript{Let $\glssymbol{frequency_vector} = (\pi_a)_{a\in\Sigma}$ be the frequency vector, which means that $\pi_a\ge 0$ for all
    $a\in\Sigma$ and $\sum_{a\in\Sigma}\pi_a = 1$, of the naturally occurring amino acids.
    If we look at
    two random positions in two randomly selected proteins, the probability that we see the ordered pair
    $(a,b)$ is then $\pi_a\cdot\pi_b$.
    The entries of $\pi$ are also called \emph{background frequencies} of the amino acids.}

\defin{background_frequencies}

\defin{fixed_evolutionary_time_period}

\skript{Now assume that we can track the fate of individual amino acids over a fixed evolutionary time
period $t$.

Let $M_t(a,b)$ be the observed frequency table of homologous amino acid pairs where we
observed $a$ at the beginning of the period and $b$ at the end of the period, such that
$\sum_{a,b} M_t(a,b) = 1$.
Usually, we count substitutions and identities twice, i.e., once in each
direction.
This ensures the symmetry of $M_t$: $M_t(a,b) = M_t(b,a)$.
The reason is that we can in
fact not track amino acids during evolution; we can only observe the state of two present-day
sequences and do not know their most recent common ancestor.
The fields of molecular evolution and
phylogenetics examine more of the resulting implications.

We can find the background frequencies as the marginals of $M_t$:
$\pi_a = \sum_{b\in\Sigma} M_t(a,b)$ and $\pi_b = \sum_{a\in\Sigma} M_t(a,b)$ because of symmetry.

There are two reasons why an entry $M_t(a,b)$ can be relatively large.
The first reason is the one we
are interested in: because $a$ and $b$ are similar.
The second reason is that simply $a$ or $b$ could
be frequent amino acids.
To remove the second effect, we consider the so-called \emph{likelihood
ratio} $M_t(a,b)/(\pi_a\cdot\pi_b)$, which relates the probability of the pair $(a,b)$ in an
evolutionary model to its probability in a random model.}

\defin{likelihood_ratio}

\skript{We declare that $a$ and $b$ are similar if
this ratio exceeds $1$ and that they are dissimilar if this ratio is below $1$.}



\skript{To obtain an additive
function, we take the logarithm and define the \emph{log-odds score matrix} with respect to time
period $t$ by
\[
    \mathrm{score}_t(a,b) \coloneqq \log\!\left(\frac{M_t(a,b)}{\pi_a\cdot\pi_b}\right).
\]
The time parameter $t$ should be chosen in such a way that the score matrix is optimal for the
problem at hand: If we want to compare two sequences whose most recent common ancestor existed, say,
$530$ million years ago, then we should ideally use a score matrix constructed from sequence pairs
with distance $t \approx 1060$ million years.}

\begin{defi}[\glsfirst*{log_odds_score_matrix}]
    \glsF*{log_odds_score_matrix}

    \begin{equation}\label{eq:log_odds_score_matrix}
        \mathrm{score}_t(a,b):=
        \log\!\left(
                  \frac{M_t(a,b)}{\pi_a\cdot\pi_b}
        \right)
    \end{equation}
\end{defi}
\begin{bem}
    For convenience, scores are often scaled by a constant factor and then rounded to integers.
    The
    score unit is called a \emph{bit} if the score is obtained by a $\log_2(\cdot)$ operation.
    When
    multiplied by a factor of three, say, we get scores in \emph{third-bits}.
\end{bem}

\defin{third_bits}

\begin{bem}
    Since we cannot easily observe how DNA or proteins change over millions of years, constructing
    a score matrix is somewhat difficult.
    Since the pioneering work of Margaret Dayhoff and colleagues in
    the 1970s, Markov processes are used to model molecular evolution.
    Methods that allow to integrate
    information from sequences of different divergence times in a consistent fashion have been developed
    more recently.
\end{bem}

\begin{bem}
    Important score matrix families for proteins are the PAM($t$) (Dayhoff et al., 1978) and the
    BLOSUM($s$) (Henikoff and Henikoff, 1992) matrices.
    Here $t$ is a divergence time parameter and $s$
    is a clustering similarity parameter inversely correlated to $t$.
    (The inventors of BLOSUM have
    chosen to index their matrices in a different way.) The BLOSUM matrices are the most widely used ones
    for protein sequence comparison.
\end{bem}


\section{Non-symmetric score matrices}\label{sec:non-symmetric-score-matrices}

\skript{We usually demand that similarity functions and hence score matrices are symmetric.
In some cases,
however, there are good reasons to choose a non-symmetric similarity function.
Often then, the
unsymmetry does not stem from an unsymmetry of $M_t$ (for reasons explained above), but from
different background frequencies.
We mention an example.

Assume that we have a particular protein fragment (the ``query'') that forms a transmembrane helix.
The amino acid composition in membrane regions differs strongly from that of the ``average''
protein: It is hydrophobically biased.
Assume that we want to look for similar fragments in a large
database of peptide sequences (of unknown or ``average'' type).
It follows that the matrix $M_t$ of
pair frequencies should be one derived from an evolutionary process acting on transmembrane helices
and that two different types of background frequencies should be used: The query background
frequencies $\tau$ are those of the transmembrane model, different from the background frequencies
$\pi$ of the database.
It follows that we should use
\[
    \mathrm{score}_t(a,b) \coloneqq \log\!\left(\frac{M_t(a,b)}{\tau_a\cdot\pi_b}\right),
\]
or a rounded multiple thereof, so that now $\mathrm{score}(a,b)\ne\mathrm{score}(b,a)$ in general.
Table~4.2 shows the SLIM161 matrix for transmembrane helices (Müller et al., 2001).
The $161$ is
the time parameter $t$ referring to the expected number of evolutionary events per $100$ positions.}


\begin{mytable}
    \begin{tabular}{c|rrrrrrrrrrrrrrrrrrrr}
        & A  & R  & N  & D   & C  & Q  & E   & G  & H  & I  & L  & K   & M  & F  & P   & S  & T  & W  & Y  & V  \\
        \hline
        A & 5  & -8 & -5 & -10 & 3  & -7 & -11 & 0  & -5 & 1  & 1  & -10 & 1  & 1  & -5  & 2  & 1  & -2 & -3 & 2  \\
        R & -3 & 10 & -4 & -10 & -2 & -3 & -11 & -4 & -4 & -2 & -1 & -3  & -1 & -2 & -7  & -4 & -3 & -2 & -3 & -3 \\
        N & -1 & -5 & 8  & -2  & 0  & -2 & -7  & -2 & 2  & -1 & -1 & -6  & 0  & 1  & -5  & 2  & 0  & -2 & 2  & -2 \\
        D & -3 & -9 & 1  & 9   & -4 & -2 & 1   & -2 & -2 & -2 & -2 & -6  & -2 & -1 & -5  & -3 & -3 & -3 & -2 & -2 \\
        C & 2  & -8 & -5 & -11 & 11 & -7 & -12 & -3 & -8 & 0  & 1  & -12 & 1  & 3  & -11 & 2  & 0  & -1 & 0  & 1  \\
        Q & -1 & -2 & 0  & -3  & 0  & 7  & -4  & -2 & 0  & 0  & 0  & -4  & 2  & 2  & -4  & 0  & -1 & 4  & 1  & 0  \\
        E & -3 & -7 & -2 & 3   & -2 & -1 & 7   & -3 & -1 & -2 & -2 & -8  & -1 & 0  & -4  & -1 & -3 & -1 & -1 & -2 \\
        G & 1  & -7 & -4 & -7  & 0  & -5 & -10 & 7  & -7 & -1 & 0  & -8  & 0  & 1  & -5  & 1  & -1 & -3 & -2 & -1 \\
        H & -1 & -5 & 3  & -5  & -2 & -1 & -5  & -4 & 10 & -1 & -1 & -8  & -1 & 2  & -7  & -1 & -2 & 1  & 5  & -2 \\
        I & -1 & -9 & -7 & -11 & -1 & -7 & -12 & -4 & -7 & 6  & 3  & -11 & 4  & 2  & -6  & -3 & -1 & -2 & -3 & 4  \\
        L & -1 & -8 & -6 & -10 & 1  & -7 & -12 & -4 & -7 & 4  & 5  & -11 & 4  & 3  & -7  & -3 & -1 & -1 & -2 & 2  \\
        K & -1 & 2  & -1 & -4  & -2 & 0  & -7  & -1 & -3 & 0  & -1 & 6   & 0  & 0  & -1  & -1 & -1 & -1 & 1  & -2 \\
        M & -1 & -7 & -6 & -10 & 1  & -5 & -11 & -3 & -7 & 4  & 4  & -11 & 7  & 3  & -7  & -2 & 0  & -1 & -2 & 2  \\
        F & -1 & -9 & -5 & -10 & 2  & -6 & -11 & -3 & -4 & 1  & 2  & -11 & 2  & 8  & -7  & -2 & -2 & 3  & 4  & 0  \\
        P & -3 & -9 & -7 & -9  & -7 & -8 & -10 & -4 & -9 & -2 & -3 & -8  & -3 & -2 & 11  & -3 & -3 & -3 & -4 & -2 \\
        S & 2  & -8 & -1 & -8  & 4  & -5 & -9  & 0  & -4 & 0  & 0  & -9  & 0  & 1  & -5  & 6  & 1  & -2 & -1 & 0  \\
        T & 2  & -7 & -3 & -8  & 2  & -6 & -10 & -2 & -5 & 2  & 2  & -9  & 3  & 2  & -4  & 2  & 4  & -4 & -2 & 2  \\
        W & -3 & -7 & -6 & -10 & 0  & -2 & -11 & -5 & -3 & -1 & 0  & -10 & 0  & 4  & -6  & -4 & -5 & 15 & 2  & -2 \\
        Y & -4 & -8 & -2 & -9  & 1  & -5 & -10 & -5 & 0  & -2 & -1 & -8  & -1 & 6  & -7  & -3 & -3 & 2  & 11 & -2 \\
        V & 1  & -9 & -7 & -10 & 1  & -7 & -11 & -4 & -7 & 5  & 3  & -12 & 3  & 2  & -6  & -2 & 0  & -2 & -3 & 5  \\
    \end{tabular}
    \caption{The non-symmetric SLIM~161 score matrix. Scores are given in third-bits, i.e., as
        $\mathrm{score}(a,b) = \mathrm{round}\!\big(3\cdot\log_2(M(a,b)/(\tau_a\pi_b))\big)$.}
\end{mytable}


\section{Position-Specific Scores}\label{sec:position-specific-scores}


\skript{The match/mismatch score in all our alignment algorithms depends on a (sensible) similarity function
$\mathrm{score}$ as defined in Definition~4.1, often given by a log-odds score matrix.
However, this
does not need to be the case.
The algorithm does not change in any way if indel and match/mismatch
scores depend also on the position in the sequences.
Therefore we can replace
$\mathrm{score}(x[i], y[j])$ by $\mathrm{score}(i,j)$ and worry later about where to obtain reasonable
score values.

A useful application of this observation is as follows: Local alignment is often used in large-scale
database searches to find remote homologs of a given protein.
Transmembrane (TM) proteins have an
unusual sequence composition (which has a large influence on the score matrix entries, as pointed out
in Section~4.2).
According to Müller et al.\ (2001), one can increase the sensitivity of the
homology search by using a bipartite scoring scheme with a (non-symmetric) transmembrane helix
specific scoring matrix (such as SLIM) for the TM helices and a general-purpose scoring matrix (such
as BLOSUM) for the remaining regions of the query protein; see Figure~4.1. As a consequence, the
score of aligning $x[i]$ with $y[j]$ does not only depend on the amino acids $x[i]$ and $y[j]$
themselves, but also on the position $i$ within the query sequence $x$.}

%\begin{figure}[H]
%    \centering
%    \caption{Bipartite scoring scheme for detection of homologous transmembrane proteins due to Müller
%    et al.\ (2001). The figure represents the Smith-Waterman alignment matrix and indicates which scoring
%    matrix is used for which query positions (rows): In transmembrane helices (TM), the
%    transmembrane-specific scoring matrix SLIM is used, elsewhere the general-purpose matrix BLOSUM is
%    used.}
%\end{figure}

